% File: ai-nhap/2026-11-15_14_de_tham_khao_tan_phu_de_1_va_tan_phu_de_2.tex
% Nguồn: Tuyển tập đề tham khảo môn Toán Tuyển sinh 10 TP.HCM 2025-2026
% Gồm 14 bài: Quận Tân Phú - Đề 1 (7 bài) và Quận Tân Phú - Đề 2 (7 bài)

% ==============================================================================
% PHẦN 1: ĐỀ THAM KHẢO TUYỂN SINH 10 QUẬN TÂN PHÚ – ĐỀ 1 (7 BÀI)
% ==============================================================================

\begin{bt}%[9D6H18-2]
% nguon: Đề tham khảo Tuyển sinh 10 năm học 2025-2026 Sở GD&ĐT TP.HCM – Quận Tân Phú - Đề 1 (Bài 1)
Cho hàm số bậc hai $y = \dfrac{x^2}{2}$ có đồ thị là parabol $(P)$.
\begin{enumerate}[a)]
    \item Vẽ đồ thị $(P)$ trên mặt phẳng tọa độ $Oxy$.
    \item Tìm tọa độ các điểm thuộc $(P)$ khác gốc tọa độ $O$ và có hoành độ gấp đôi tung độ.
\end{enumerate}
\loigiai{
\begin{enumerate}[a)]
    \item Bảng giá trị tương ứng của hàm số $y = \dfrac{x^2}{2}$:
    \begin{center}
    \begin{tabular}{|c|c|c|c|c|c|}
        \hline
        $x$ & $-4$ & $-2$ & $0$ & $2$ & $4$ \\
        \hline
        $y = \dfrac{x^2}{2}$ & $8$ & $2$ & $0$ & $2$ & $8$ \\
        \hline
    \end{tabular}
    \end{center}
    Đồ thị $(P)$ là một parabol có đỉnh $O(0; 0)$, đối xứng qua trục tung $Oy$ và bề lõm quay lên trên.
    \item Điểm $M(x; y)$ thuộc $(P)$ có hoành độ gấp đôi tung độ khi và chỉ khi:
    $x = 2y \iff y = \dfrac{x}{2}$.
    Thay vào phương trình của $(P)$, ta có:
    $\dfrac{x^2}{2} = \dfrac{x}{2} \iff x^2 - x = 0 \iff x(x - 1) = 0 \iff \left[\begin{aligned} &x = 0 \\ &x = 1. \end{aligned}\right.$
    Vì điểm cần tìm khác gốc tọa độ $O(0; 0)$ nên $x = 1 \implies y = \dfrac{1^2}{2} = \dfrac{1}{2}$.
    Vậy điểm thỏa mãn yêu cầu là $M\left(1; \dfrac{1}{2}\right)$.
\end{enumerate}
}
\end{bt}

\begin{bt}%[9D6H20-2]
% nguon: Đề tham khảo Tuyển sinh 10 năm học 2025-2026 Sở GD&ĐT TP.HCM – Quận Tân Phú - Đề 1 (Bài 2)
Cho phương trình bậc hai $2x^2 - 5x - 3 = 0$.
\begin{enumerate}[a)]
    \item Chứng minh phương trình luôn có hai nghiệm phân biệt $x_1, x_2$.
    \item Không giải phương trình, hãy tính giá trị của biểu thức $A = \dfrac{x_1 - 2}{x_2} + \dfrac{x_2 - 2}{x_1}$.
\end{enumerate}
\loigiai{
\begin{enumerate}[a)]
    \item Phương trình $2x^2 - 5x - 3 = 0$ có các hệ số $a = 2$, $c = -3$.
    Vì $a \cdot c = 2 \cdot (-3) = -6 < 0$ nên phương trình luôn có hai nghiệm phân biệt trái dấu $x_1, x_2$.
    \item Áp dụng hệ thức Viète cho phương trình $2x^2 - 5x - 3 = 0$, ta có:
    $\begin{cases} S = x_1 + x_2 = \dfrac{5}{2} \\ P = x_1 x_2 = -\dfrac{3}{2}. \end{cases}$
    Quy đồng mẫu thức biểu thức $A$, ta được:
    $A = \dfrac{x_1(x_1 - 2) + x_2(x_2 - 2)}{x_1 x_2} = \dfrac{x_1^2 + x_2^2 - 2(x_1 + x_2)}{x_1 x_2} = \dfrac{(x_1 + x_2)^2 - 2x_1 x_2 - 2(x_1 + x_2)}{x_1 x_2}$.
    Thay $S = \dfrac{5}{2}$ và $P = -\dfrac{3}{2}$ vào tử số:
    $S^2 - 2P - 2S = \left(\dfrac{5}{2}\right)^2 - 2 \cdot \left(-\dfrac{3}{2}\right) - 2 \cdot \dfrac{5}{2} = \dfrac{25}{4} + 3 - 5 = \dfrac{25}{4} - 2 = \dfrac{17}{4}$.
    Do đó giá trị biểu thức $A$ là:
    $A = \dfrac{\frac{17}{4}}{-\frac{3}{2}} = -\dfrac{17}{4} \cdot \dfrac{2}{3} = -\dfrac{17}{6}$.
\end{enumerate}
}
\end{bt}

\begin{bt}%[9D8H26-3]
% nguon: Đề tham khảo Tuyển sinh 10 năm học 2025-2026 Sở GD&ĐT TP.HCM – Quận Tân Phú - Đề 1 (Bài 3)
Hai bạn An và Bình cùng chơi trò chơi phóng phi tiêu tính điểm trên một bảng tiêu chuẩn gồm $20$ cung tròn ứng với các điểm số từ $1$ đến $20$. Quy tắc tính điểm:
\begin{itemize}
    \item Vòng hồng tâm (Bull's eye): $50$ điểm.
    \item Vòng rìa hồng tâm (Bull): $25$ điểm.
    \item Vòng nhân đôi (Double ring): nhận số điểm gấp đôi con số của cung tròn đó.
    \item Vòng nhân ba (Triple ring): nhận số điểm gấp ba con số của cung tròn đó.
\end{itemize}
Bạn An phóng ba lần: lần 1 vào vòng Triple ring của cung $4$ điểm, lần 2 vào vòng Triple ring của cung $12$ điểm, lần 3 vào vòng Bull.
Bạn Bình phóng ba lần: lần 1 vào vòng Triple ring của cung $9$ điểm, lần 2 vào vòng Double ring của cung $7$ điểm, lần 3 vào vòng Bull.
Hỏi sau ba lần ném, bạn nào có tổng số điểm cao hơn?
\loigiai{
Tổng số điểm của bạn An sau ba lần phóng là:
$S_{\text{An}} = 3 \cdot 4 + 3 \cdot 12 + 25 = 12 + 36 + 25 = 73$ (điểm).
Tổng số điểm của bạn Bình sau ba lần phóng là:
$S_{\text{Bình}} = 3 \cdot 9 + 2 \cdot 7 + 25 = 27 + 14 + 25 = 66$ (điểm).
Vì $73 > 66$ nên sau ba lần ném bạn An có tổng số điểm cao hơn bạn Bình.
}
\end{bt}

\begin{bt}%[9D1H3-4]
% nguon: Đề tham khảo Tuyển sinh 10 năm học 2025-2026 Sở GD&ĐT TP.HCM – Quận Tân Phú - Đề 1 (Bài 4)
Một khu vui chơi đưa ra hai gói giá vé ưu đãi cho các nhóm gia đình như sau:
\begin{itemize}
    \item Chương trình 1: Giá vé người lớn là $100\,000$ đồng/vé; mỗi vé trẻ em có giá gốc $50\,000$ đồng được giảm $20\,000$ đồng (còn $30\,000$ đồng/vé).
    \item Chương trình 2: Giá vé người lớn là $80\,000$ đồng/vé; giá vé trẻ em giữ nguyên là $50\,000$ đồng/vé.
\end{itemize}
\begin{enumerate}[a)]
    \item Gọi $x$ là số người lớn trong nhóm ($x \in \mathbb{N}^*$). Hãy lập biểu thức biểu thị tổng số tiền vé nhóm phải trả theo Chương trình 1, biết rằng nhóm có $4$ trẻ em.
    \item Hỏi nhóm đó cần có tối thiểu bao nhiêu người lớn để việc lựa chọn mua vé theo Chương trình 2 tiết kiệm chi phí hơn Chương trình 1?
\end{enumerate}
\loigiai{
\begin{enumerate}[a)]
    \item Theo Chương trình 1, giá vé của mỗi trẻ em là:
    $50\,000 - 20\,000 = 30\,000$ (đồng).
    Tổng số tiền vé nhóm gồm $x$ người lớn và $4$ trẻ em phải trả theo Chương trình 1 là:
    $T_1(x) = 100\,000x + 4 \cdot 30\,000 = 100\,000x + 120\,000$ (đồng).
    \item Theo Chương trình 2, tổng số tiền vé nhóm phải trả là:
    $T_2(x) = 80\,000x + 4 \cdot 50\,000 = 80\,000x + 200\,000$ (đồng).
    Để Chương trình 2 có lợi hơn (chi phí thấp hơn) Chương trình 1 thì:
    $T_2(x) < T_1(x) \iff 80\,000x + 200\,000 < 100\,000x + 120\,000 \iff 20\,000x > 80\,000 \iff x > 4$.
    Vì $x$ là số nguyên dương nên giá trị nhỏ nhất của $x$ là $5$.
    Vậy nhóm gia đình cần có tối thiểu $5$ người lớn để chọn Chương trình 2 có lợi hơn.
\end{enumerate}
}
\end{bt}

\begin{bt}%[9D1H3-4]
% nguon: Đề tham khảo Tuyển sinh 10 năm học 2025-2026 Sở GD&ĐT TP.HCM – Quận Tân Phú - Đề 1 (Bài 5)
Theo kế hoạch, hai tổ sản xuất được giao nhiệm vụ hoàn thành tổng cộng $700$ sản phẩm trong một thời gian nhất định. Nhờ áp dụng cải tiến kỹ thuật, tổ I đã sản xuất vượt mức $12\%$ so với kế hoạch, trong khi tổ II gặp khó khăn nên sản lượng giảm $5\%$ so với kế hoạch. Tuy nhiên, tính chung cả hai tổ vẫn hoàn thành vượt mức $16$ sản phẩm. Hãy tính số sản phẩm làm vượt mức của tổ I và số sản phẩm bị hụt giảm của tổ II so với chỉ tiêu kế hoạch ban đầu.
\loigiai{
Gọi số sản phẩm được giao theo kế hoạch của tổ I là $x$ và của tổ II là $y$ ($x, y \in \mathbb{N}^*$, $x, y < 700$).
Vì tổng số sản phẩm theo kế hoạch của hai tổ là $700$ nên ta có phương trình:
$x + y = 700$. \quad (1)
Số sản phẩm tổ I làm vượt mức là: $12\% \cdot x = 0{,}12x$ (sản phẩm).
Số sản phẩm tổ II bị giảm là: $5\% \cdot y = 0{,}05y$ (sản phẩm).
Do cả hai tổ hoàn thành vượt mức $16$ sản phẩm so với kế hoạch nên:
$0{,}12x - 0{,}05y = 16$. \quad (2)
Từ (1) và (2) ta có hệ phương trình:
$\begin{cases} x + y = 700 \\ 0{,}12x - 0{,}05y = 16 \end{cases} \iff \begin{cases} 0{,}05x + 0{,}05y = 35 \\ 0{,}12x - 0{,}05y = 16 \end{cases} \iff \begin{cases} 0{,}17x = 51 \\ y = 700 - x \end{cases} \iff \begin{cases} x = 300 \\ y = 400 \end{cases}$ (thỏa mãn điều kiện).
Số sản phẩm tổ I làm vượt mức là: $0{,}12 \cdot 300 = 36$ (sản phẩm).
Số sản phẩm tổ II bị giảm hụt là: $0{,}05 \cdot 400 = 20$ (sản phẩm).
}
\end{bt}

\begin{bt}%[9H0V31-3]
% nguon: Đề tham khảo Tuyển sinh 10 năm học 2025-2026 Sở GD&ĐT TP.HCM – Quận Tân Phú - Đề 1 (Bài 6)
Kim tự tháp kính tại bảo tàng Louvre (Paris, Pháp) có dạng một hình chóp tứ giác đều với đáy là hình vuông cạnh dài $34\text{ m}$ và chiều cao của kim tự tháp là $21{,}3\text{ m}$. Hãy tính thể tích không gian bên trong của kim tự tháp này.
\loigiai{
Diện tích mặt đáy hình vuông của kim tự tháp là:
$S_{\text{đáy}} = 34^2 = 1156\text{ (m}^2\text{)}$.
Thể tích không gian bên trong kim tự tháp hình chóp tứ giác đều là:
$V = \dfrac{1}{3} S_{\text{đáy}} \cdot h = \dfrac{1}{3} \cdot 1156 \cdot 21{,}3 = 1156 \cdot 7{,}1 = 8207{,}6\text{ m}^3$.
Vậy thể tích của kim tự tháp Louvre là $8207{,}6\text{ m}^3$.
}
\end{bt}

\begin{bt}%[9H9C29-3]
% nguon: Đề tham khảo Tuyển sinh 10 năm học 2025-2026 Sở GD&ĐT TP.HCM – Quận Tân Phú - Đề 1 (Bài 7)
Cho tam giác $ABC$ nhọn nội tiếp đường tròn $(O; R)$ ($AB > AC$) có đường kính $AM$. Hai đường cao $CE$ và $AD$ của tam giác cắt nhau tại trực tâm $H$.
\begin{enumerate}[a)]
    \item Chứng minh tứ giác $BEHD$ là tứ giác nội tiếp và $\triangle AEH \backsim \triangle ACM$.
    \item Gọi $I$ là trung điểm của đoạn thẳng $HM$ và $HM$ cắt đường tròn $(O)$ tại điểm $N$. Chứng minh ba điểm $B, I, C$ thẳng hàng và $BC^2 = 4IH \cdot IN$.
    \item Giả sử $\widehat{BAC} = 60^\circ$. Chứng minh độ dài đoạn thẳng $AH = R$.
\end{enumerate}
\loigiai{
\begin{enumerate}[a)]
    \item Vì $AD \perp BC$ tại $D$ và $CE \perp AB$ tại $E$ nên $\widehat{BDH} = 90^\circ$ và $\widehat{BEH} = 90^\circ$.
    Tứ giác $BEHD$ có $\widehat{BDH} + \widehat{BEH} = 90^\circ + 90^\circ = 180^\circ$ nên là tứ giác nội tiếp.
    Lại có $AM$ là đường kính của $(O)$ nên tam giác $ACM$ nội tiếp chắn nửa đường tròn, suy ra $\widehat{ACM} = 90^\circ$.
    Xét $\triangle AEH$ vuông tại $E$ và $\triangle ACM$ vuông tại $C$, ta có:
    $\widehat{AHE} = \widehat{AMC}$ (cùng phụ với góc $\widehat{BAD}$ do tứ giác $ABMC$ nội tiếp).
    Do đó $\triangle AEH \backsim \triangle ACM$ (g.g).
    \item Ta có $BH \perp AC$ (do $H$ là trực tâm $\triangle ABC$) và $MC \perp AC$ (do $\widehat{ACM} = 90^\circ$), suy ra $BH \parallel MC$.
    Tương tự $CH \perp AB$ và $MB \perp AB$ (do $\widehat{ABM} = 90^\circ$), suy ra $CH \parallel MB$.
    Do đó tứ giác $BHCM$ là hình bình hành.
    Vì $I$ là trung điểm của đường chéo $HM$ nên $I$ cũng là trung điểm của đường chéo $BC$, suy ra ba điểm $B, I, C$ thẳng hàng và $BC = 2IB = 2IC$.
    Áp dụng hệ thức phương tích của điểm $I$ đối với đường tròn $(O)$, ta có hai dây cung $BC$ và $MN$ cắt nhau tại $I$ nên:
    $IB \cdot IC = IM \cdot IN$.
    Mà $IB = IC = \dfrac{BC}{2}$ và $IM = IH$ (do $I$ là trung điểm $HM$), suy ra:
    $\left(\dfrac{BC}{2}\right)^2 = IH \cdot IN \iff \dfrac{BC^2}{4} = IH \cdot IN \iff BC^2 = 4IH \cdot IN$ (đpcm).
    \item Khi $\widehat{BAC} = 60^\circ$, xét tam giác $OBC$ cân tại $O$ có góc ở tâm $\widehat{BOC} = 2\widehat{BAC} = 120^\circ$.
    Đoạn thẳng $OI \perp BC$ tại trung điểm $I$, ta có:
    $OI = OB \cdot \cos 60^\circ = R \cdot \dfrac{1}{2} = \dfrac{R}{2}$.
    Mặt khác, trong hình bình hành $BHCM$, đoạn nối trung điểm $OI$ là đường trung bình của tam giác $AHM$ (vì $O$ là trung điểm $AM$ và $I$ là trung điểm $HM$), suy ra:
    $AH = 2OI = 2 \cdot \dfrac{R}{2} = R$ (đpcm).
\end{enumerate}
}
\end{bt}


% ==============================================================================
% PHẦN 2: ĐỀ THAM KHẢO TUYỂN SINH 10 QUẬN TÂN PHÚ – ĐỀ 2 (7 BÀI)
% ==============================================================================

\begin{bt}%[9D6H18-2]
% nguon: Đề tham khảo Tuyển sinh 10 năm học 2025-2026 Sở GD&ĐT TP.HCM – Quận Tân Phú - Đề 2 (Bài 1)
Cho hàm số bậc hai $y = \dfrac{1}{4}x^2$ có đồ thị là đường parabol $(P)$.
\begin{enumerate}[a)]
    \item Vẽ đồ thị $(P)$ trên mặt phẳng tọa độ $Oxy$.
    \item Tìm tọa độ điểm $A$ thuộc đồ thị $(P)$ sao cho hoành độ của điểm $A$ lớn hơn tung độ của nó đúng một đơn vị.
\end{enumerate}
\loigiai{
\begin{enumerate}[a)]
    \item Bảng giá trị tương ứng của hàm số $y = \dfrac{1}{4}x^2$:
    \begin{center}
    \begin{tabular}{|c|c|c|c|c|c|}
        \hline
        $x$ & $-4$ & $-2$ & $0$ & $2$ & $4$ \\
        \hline
        $y = \dfrac{1}{4}x^2$ & $4$ & $1$ & $0$ & $1$ & $4$ \\
        \hline
    \end{tabular}
    \end{center}
    Đồ thị $(P)$ là một parabol có đỉnh tại gốc tọa độ $O(0; 0)$, đối xứng qua trục tung $Oy$ và quay bề lõm lên phía trên.
    \item Điểm $A(x; y)$ thuộc $(P)$ có hoành độ lớn hơn tung độ một đơn vị thỏa mãn:
    $x - y = 1 \iff y = x - 1$.
    Thay vào phương trình của $(P)$, ta có:
    $\dfrac{1}{4}x^2 = x - 1 \iff x^2 = 4x - 4 \iff x^2 - 4x + 4 = 0 \iff (x - 2)^2 = 0 \iff x = 2$.
    Với $x = 2 \implies y = 2 - 1 = 1$.
    Vậy điểm cần tìm là $A(2; 1)$.
\end{enumerate}
}
\end{bt}

\begin{bt}%[9D6H20-2]
% nguon: Đề tham khảo Tuyển sinh 10 năm học 2025-2026 Sở GD&ĐT TP.HCM – Quận Tân Phú - Đề 2 (Bài 2)
Cho phương trình bậc hai $2x^2 - 4x - 1 = 0$ có hai nghiệm phân biệt $x_1, x_2$. Không giải phương trình, hãy tính giá trị của biểu thức $T = \dfrac{x_1}{x_2} + \dfrac{x_2}{x_1} + \dfrac{5}{2}(x_1 - x_2)^2$.
\loigiai{
Phương trình $2x^2 - 4x - 1 = 0$ có $a \cdot c = 2 \cdot (-1) = -2 < 0$ nên luôn có hai nghiệm phân biệt trái dấu $x_1, x_2$.
Theo hệ thức Viète, ta có:
$\begin{cases} S = x_1 + x_2 = -\dfrac{b}{a} = \dfrac{4}{2} = 2 \\ P = x_1 x_2 = \dfrac{c}{a} = -\dfrac{1}{2}. \end{cases}$
Biến đổi các số hạng trong biểu thức $T$:
$\dfrac{x_1}{x_2} + \dfrac{x_2}{x_1} = \dfrac{x_1^2 + x_2^2}{x_1 x_2} = \dfrac{(x_1 + x_2)^2 - 2x_1 x_2}{x_1 x_2} = \dfrac{S^2 - 2P}{P}$.
Thay $S = 2$ và $P = -\dfrac{1}{2}$ vào:
$\dfrac{S^2 - 2P}{P} = \dfrac{2^2 - 2 \cdot \left(-\dfrac{1}{2}\right)}{-\dfrac{1}{2}} = \dfrac{4 + 1}{-\dfrac{1}{2}} = \dfrac{5}{-\dfrac{1}{2}} = -10$.
Số hạng thứ hai:
$(x_1 - x_2)^2 = (x_1 + x_2)^2 - 4x_1 x_2 = S^2 - 4P = 2^2 - 4 \cdot \left(-\dfrac{1}{2}\right) = 4 + 2 = 6$.
Do đó:
$\dfrac{5}{2}(x_1 - x_2)^2 = \dfrac{5}{2} \cdot 6 = 15$.
Vậy giá trị của biểu thức $T$ là:
$T = -10 + 15 = 5$.
}
\end{bt}

\begin{bt}%[9D8H26-3]
% nguon: Đề tham khảo Tuyển sinh 10 năm học 2025-2026 Sở GD&ĐT TP.HCM – Quận Tân Phú - Đề 2 (Bài 3)
Số lượng học sinh tham gia giải thể thao cấp trường của một trường THCS được thống kê theo từng khối lớp như sau: khối 6 có $9$ nam và $8$ nữ; khối 7 có $8$ nam và $7$ nữ; khối 8 có $7$ nam và $5$ nữ; khối 9 có $5$ nam và $4$ nữ. Chọn ngẫu nhiên một học sinh trong danh sách tham gia giải thể thao của trường. Tính xác suất của các biến cố sau:
\begin{enumerate}[a)]
    \item Biến cố $A$: ``Học sinh được chọn là nam''.
    \item Biến cố $B$: ``Học sinh được chọn thuộc khối 6''.
    \item Biến cố $C$: ``Học sinh được chọn là nữ và không thuộc khối 9''.
\end{enumerate}
\loigiai{
Tổng số học sinh tham gia thi đấu thể thao của toàn trường là:
$n(\Omega) = (9 + 8) + (8 + 7) + (7 + 5) + (5 + 4) = 17 + 15 + 12 + 9 = 53$ (học sinh).
\begin{enumerate}[a)]
    \item Tổng số học sinh nam tham gia là:
    $n(A) = 9 + 8 + 7 + 5 = 29$ (học sinh).
    Xác suất của biến cố $A$ là:
    $P(A) = \dfrac{n(A)}{n(\Omega)} = \dfrac{29}{53}$.
    \item Tổng số học sinh thuộc khối 6 tham gia là:
    $n(B) = 9 + 8 = 17$ (học sinh).
    Xác suất của biến cố $B$ là:
    $P(B) = \dfrac{n(B)}{n(\Omega)} = \dfrac{17}{53}$.
    \item Học sinh là nữ và không thuộc khối 9 gồm học sinh nữ của các khối 6, 7 và 8.
    Số kết quả thuận lợi cho biến cố $C$ là:
    $n(C) = 8 + 7 + 5 = 20$ (học sinh).
    Xác suất của biến cố $C$ là:
    $P(C) = \dfrac{n(C)}{n(\Omega)} = \dfrac{20}{53}$.
\end{enumerate}
}
\end{bt}

\begin{bt}%[9H0V31-3]
% nguon: Đề tham khảo Tuyển sinh 10 năm học 2025-2026 Sở GD&ĐT TP.HCM – Quận Tân Phú - Đề 2 (Bài 4)
Từ một khối gỗ có dạng hình lập phương cạnh $x$ (cm), người thợ tiện khoét bỏ đi một phần ruột có dạng hình nón có đỉnh chạm vào tâm mặt đáy dưới của khối gỗ và đáy hình nón nằm trọn trên mặt phẳng đáy trên với đường kính là $2y$ (cm) ($x \ge 2y$).
\begin{enumerate}[a)]
    \item Viết biểu thức biểu thị thể tích $V$ của phần khối gỗ còn lại theo $x$ và $y$.
    \item Biết rằng thể tích khối gỗ hình lập phương ban đầu bằng $27\,000\text{ cm}^3$ và thể tích phần gỗ còn lại sau khi khoét xấp xỉ $6158\text{ cm}^3$. Tính độ dài cạnh $x$ và bán kính đáy $y$ của hình nón (kết quả làm tròn đến hàng phần mười). Lấy $\pi \approx 3{,}1416$.
\end{enumerate}
\loigiai{
\begin{enumerate}[a)]
    \item Thể tích của khối gỗ hình lập phương cạnh $x$ là: $V_1 = x^3\text{ (cm}^3\text{)}$.
    Hình nón có chiều cao bằng chiều cao khối gỗ ($h = x$) và bán kính đáy là $y$.
    Thể tích khối gỗ hình nón bị khoét bỏ là:
    $V_2 = \dfrac{1}{3}\pi y^2 x\text{ (cm}^3\text{)}$.
    Thể tích phần khối gỗ còn lại là:
    $V = V_1 - V_2 = x^3 - \dfrac{1}{3}\pi y^2 x\text{ (cm}^3\text{)}$.
    \item Theo đề bài, thể tích khối lập phương ban đầu là $27\,000\text{ cm}^3$:
    $x^3 = 27\,000 \iff x = 30\text{ cm}$.
    Thể tích phần gỗ còn lại là $6158\text{ cm}^3$, ta có phương trình:
    $27\,000 - \dfrac{1}{3}\pi y^2 \cdot 30 = 6158 \iff 10\pi y^2 = 27\,000 - 6158 \iff 10\pi y^2 = 20\,842$.
    Suy ra:
    $y^2 = \dfrac{20\,842}{10\pi} \approx \dfrac{20\,842}{31{,}416} \approx 663{,}42 \implies y = \sqrt{663{,}42} \approx 25{,}8\text{ cm}$.
    Vậy $x = 30\text{ cm}$ và $y \approx 25{,}8\text{ cm}$.
\end{enumerate}
}
\end{bt}

\begin{bt}%[9H0V31-3]
% nguon: Đề tham khảo Tuyển sinh 10 năm học 2025-2026 Sở GD&ĐT TP.HCM – Quận Tân Phú - Đề 2 (Bài 5)
Một quả cầu đặc có bán kính $R = 7\text{ cm}$ nằm khít bên trong một chiếc bình dạng hình trụ có chiều cao $h = 2R = 14\text{ cm}$ và bán kính đường tròn đáy bằng $R$. Người ta đổ nước vào bình cho đến khi ngập kín quả cầu và đầy tràn miệng bình. Lấy $\pi \approx 3{,}1416$.
\begin{enumerate}[a)]
    \item Tính thể tích của chiếc bình hình trụ (kết quả làm tròn đến hàng phần trăm của $\text{cm}^3$).
    \item Tính thể tích lượng nước thực tế đang chứa bên trong chiếc bình (kết quả làm tròn đến hàng phần trăm của $\text{cm}^3$).
\end{enumerate}
\loigiai{
\begin{enumerate}[a)]
    \item Chiều cao hình trụ là $h = 2R = 2 \cdot 7 = 14\text{ cm}$.
    Thể tích của chiếc bình hình trụ là:
    $V_{\text{bình}} = \pi R^2 h = \pi \cdot 7^2 \cdot 14 = 686\pi \approx 2155{,}13\text{ cm}^3$.
    \item Thể tích của quả cầu đặc bên trong bình là:
    $V_{\text{cầu}} = \dfrac{4}{3}\pi R^3 = \dfrac{4}{3}\pi \cdot 7^3 = \dfrac{1372\pi}{3} \approx 1436{,}75\text{ cm}^3$.
    Lượng nước chứa trong bình bằng hiệu thể tích giữa bình hình trụ và quả cầu:
    $V_{\text{nước}} = V_{\text{bình}} - V_{\text{cầu}} = 686\pi - \dfrac{1372\pi}{3} = \dfrac{686\pi}{3} \approx 718{,}38\text{ cm}^3$.
\end{enumerate}
}
\end{bt}

\begin{bt}%[9D1H3-4]
% nguon: Đề tham khảo Tuyển sinh 10 năm học 2025-2026 Sở GD&ĐT TP.HCM – Quận Tân Phú - Đề 2 (Bài 6)
Trong một kỳ thi thử tốt nghiệp THCS gồm $50$ câu hỏi trắc nghiệm, mỗi câu trả lời đúng được cộng $20$ điểm, còn mỗi câu trả lời sai bị trừ đi $5$ điểm. Bạn Na đã tham gia làm trọn vẹn cả $50$ câu hỏi và đạt được tổng số điểm là $550$ điểm. Hỏi bạn Na đã trả lời đúng bao nhiêu câu và trả lời sai bao nhiêu câu?
\loigiai{
Gọi số câu trả lời đúng của Na là $x$ và số câu trả lời sai là $y$ ($x, y \in \mathbb{N}$, $x, y \le 50$).
Vì Na làm trọn vẹn $50$ câu hỏi nên ta có phương trình:
$x + y = 50$. \quad (1)
Mỗi câu đúng được $+20$ điểm và mỗi câu sai bị $-5$ điểm, tổng số điểm đạt được là $550$ điểm nên:
$20x - 5y = 550 \iff 4x - y = 110$. \quad (2)
Cộng hai phương trình (1) và (2) theo từng vế, ta được:
$5x = 160 \iff x = 32$.
Thay $x = 32$ vào (1), ta được:
$y = 50 - 32 = 18$ (thỏa mãn điều kiện).
Vậy bạn Na đã trả lời đúng $32$ câu và trả lời sai $18$ câu.
}
\end{bt}

\begin{bt}%[9H9C29-3]
% nguon: Đề tham khảo Tuyển sinh 10 năm học 2025-2026 Sở GD&ĐT TP.HCM – Quận Tân Phú - Đề 2 (Bài 7)
Từ điểm $A$ nằm bên ngoài đường tròn $(O; R)$, kẻ hai tiếp tuyến $AB, AC$ với đường tròn $(O)$ ($B, C$ là các tiếp điểm). Đoạn thẳng $AO$ cắt dây cung $BC$ tại điểm $K$.
\begin{enumerate}[a)]
    \item Chứng minh $OA \perp BC$ tại $K$ và tứ giác $ABOC$ là tứ giác nội tiếp.
    \item Kẻ đường kính $BE$ của đường tròn $(O)$, tia $AE$ cắt $(O)$ tại điểm thứ hai là $D$. Chứng minh $AB^2 = AK \cdot AO$ và $\widehat{AEO} = \widehat{AKD}$.
    \item Kẻ đường thẳng $CH$ vuông góc với đường kính $BE$ tại điểm $H$, đoạn thẳng $AE$ cắt $CH$ tại điểm $I$, tia $EC$ cắt tia $BA$ tại điểm $N$. Chứng minh $A$ là trung điểm của đoạn thẳng $BN$ và $IH = IC$.
\end{enumerate}
\loigiai{
\begin{enumerate}[a)]
    \item Vì $AB, AC$ là hai tiếp tuyến của $(O)$ nên $AB = AC$ và $OB = OC = R$.
    Do đó $OA$ là đường trung trực của $BC$, suy ra $OA \perp BC$ tại $K$.
    Lại có $\widehat{ABO} = \widehat{ACO} = 90^\circ$ nên tứ giác $ABOC$ có $\widehat{ABO} + \widehat{ACO} = 180^\circ$, do đó $ABOC$ là tứ giác nội tiếp đường tròn đường kính $AO$.
    \item Tam giác $ABO$ vuông tại $B$ có đường cao $BK$, áp dụng hệ thức lượng ta có:
    $AB^2 = AK \cdot AO$.
    Mặt khác, trong tam giác vuông $ABE$ tại $B$ (vì $BE$ là đường kính và $AB$ là tiếp tuyến) có $BD \perp AE$ (do $\widehat{BDE} = 90^\circ$), ta có $AB^2 = AD \cdot AE$.
    Từ đó suy ra $AD \cdot AE = AK \cdot AO \implies \dfrac{AD}{AO} = \dfrac{AK}{AE}$.
    Xét $\triangle ADK$ và $\triangle AOE$, ta có góc $\widehat{OAE}$ chung và $\dfrac{AD}{AO} = \dfrac{AK}{AE}$, suy ra $\triangle ADK \backsim \triangle AOE$ (c.g.c).
    Do đó $\widehat{AKD} = \widehat{AEO}$ (đpcm).
    \item Ta có $\widehat{BCE} = 90^\circ$ (góc nội tiếp chắn nửa đường tròn), suy ra $BC \perp NE$.
    Trong tam giác vuông $NBE$ tại $B$, $BC$ là đường cao hạ từ $B$, do đó $\widehat{NBC} = \widehat{NEB}$.
    Lại có $AB = AC$ nên tam giác $ABC$ cân tại $A$, từ các quan hệ góc suy ra $A$ là trung điểm của cạnh huyền $BN$ trong tam giác vuông $NBE$.
    Vì $CH \parallel BN$ (cùng vuông góc với $BE$), áp dụng định lí Thales và hệ quả vào các tam giác $EBN$ và $EAN$:
    $\dfrac{CI}{AN} = \dfrac{EI}{EA} = \dfrac{IH}{AB}$.
    Vì $AN = AB$ (do $A$ là trung điểm $BN$) nên $IC = IH$.
    Vậy $IH = IC$ (đpcm).
\end{enumerate}
}
\end{bt}
